
%!TEX root = ../main.tex

\chapter{Stand der Technik und Forschung}
\label{chap:StandTechnikUForschung}
Die Entwicklung von \ac{AR}-Brillen ist ein aktives Forschungs- und Industriefeld, in dem kommerzielle Produkte und wissenschaftliche Arbeiten gleichermaßen den technologischen Fortschritt vorantreiben. Das folgende Kapitel gibt einen Überblick über repräsentative kommerzielle \ac{AR}-Systeme sowie relevante akademische Forschungsansätze, um das in dieser Arbeit entwickelte System im Kontext des aktuellen Stands der Technik einzuordnen.

\section{Kommerzielle AR-Systeme}
\label{sec:KommerzielleSysteme}
Die kommerzielle Entwicklung von \ac{AR}-Brillen hat in den vergangenen Jahren erheblich an Dynamik gewonnen, wobei sich zwei grundlegende Produktkategorien herausgebildet haben: hochintegrierte Mixed-Reality-Headset für den professionellen Einsatz sowie leichte Consumer-Brillen mit eingeschränktem Funktionsumfang. Der Markt für \ac{NED}s wurde 2025 auf ca. 3,8 Milliarden USD geschätzt und soll bis 2035 auf 32 Milliarden USD anwachsen. Die nachfolgend beschriebenen Systeme repräsentieren den aktuellen Stand der Technik und dienen als Referenzrahmen für die Einordnung des im Rahmen dieser Arbeit entwickelten Systems.

\subsection{Microsoft HoloLens 2}
\label{subsec:HoloLens}
Die Microsoft HoloLens 2 ist ein eigenständiges Mixed-Reality-Headset, das 2019 vorgestellt wurde und primär für industrielle und professionelle Anwendungen konzipiert ist. Als Optiksystem kommen diffractive Wellenleitertechnologie (Waveguide) mit \ac{LCoS}-Lichtmaschinen zum Einsatz, die ein diagonales \ac{FoV} von 52\textdegree erreichen. Die Sensorbasis umfasst vier sichtbare Licht-Kameras für das Head-Tracking, eine Time-of-Flight-Tiefenkamera sowie zwei Infrarotkameras für Eye-Tracking, was eine präzise räumliche Interaktion ermöglicht. Mit einem Gewicht von über 500g und einem Verkaufspreis von mehreren Tausend Euro richtet sich das System explizit an Enterprise-Anwendungen in den Bereichen Industrie 4.0, Remote-Wartung und medizinische Ausbildung.

\subsection{Apple Vision Pro}
\label{subsec:AppleVision}
Die Apple Vision Pro, seit Februar 2024 auf dem Markt erhältlich, setzt auf eine Micro-OLED-Displaytechnologie mit jeweils 23 Millionen Pixeln pro Auge bei einem Pixelabstand von 7,5 Mikrometern. Also Optiksystem kommen katadioptische Pancake-Linsen zum Einsatz, die eine kompaktere Bautiefe als klassische Freiraumoptiken ermöglichen. Die Besonderheit der Vision Pro liegt in ihrem \ac{VST} Ansatz: Da das Display nicht optisch transparent ist, wird die Außenwelt über externe Kameras in Echtzeit auf die internen Displays gestreamt, was technisch kein \ac{OST} darstellt, sondern eine elektronische Umgebungswiedergabe. Das System wiegt ca. 600g und kostet ab einigen Tausend Euro, womit es sich trotz Consumer-Positionierung im professionelen Preissegment befindet.

\subsection{XREAL One}
\label{subsec:Xreal}
Die XREAL One (ehemalig Nreal) repräsentiert eine leichtere Consumer-Kategorie und setzt das Birdbath-Optikdesign ein, das einen 50\textdegree-\ac{FoV} bei einem Gesamtgewicht von unter 100g ermöglicht. Als Besonderheit verfügt das System über den selbstentwickelten X1-Spatial-Computing-Chip, der eine Motion-to-Photon-Latenz von 3ms bei 120Hz erreicht und damit eine stabile räumliche Bildverankerung ohne Ghosting ermöglicht. Die Brille verdeutlicht, dass das Birdbath-Design trotz siner optischen Einschränkungen - insbesondere der eingeschränkten Transparenz von ca 15-25\% - durch Kostenvorteile und Bildqualität relevant für Consumer-Produkte bleibt.

\subsection{Meta Ray-Ban Display}
\label{subsec:Ray-Ban}
Die 2025 vorgestellten Meta Ray-Ban Display Glasses integrieren ein monokulares In-Lens-Display mit 600x600 Pixeln, 20\textdegree \ac{FoV} und einer Spitzenhelligkeit von 5.000Nits in Brillenframes mit einem Gewicht von 69g. Das System repräsentiert eine radikal andere Design-Philosophie gegenüber HoloLens und Vision Pro: An Stelle eines großen \ac{FoV} und räumlicher Interaktion steht die Form-Faktor-Priorität, bei der das Display als diskretes Informationsdisplay fungiert. Bemerkenswert ist, dass der Lichtaustritt aus dem Display nach außen auf 2\% begrenzt ist, was eine Nutzung ohne für Dritte sichtbare Inhalte ermöglicht.

\section{Aktuelle Forschungsansätze}
\label{sec:aktuelleforschungsansätze}
Die akademische Forschung im Bereich der \ac{AR}-Displays ist ein dynamisches Feld, das auf mehreren parallel verlaufenden Fronten fundamentale Durchbrüche anstrebt. Im Folgenden werden vier der derzeit aktivsten und bedeutsamsten Forschungsrichtungen beschrieben.

\subsection{Holografische Displays}
\label{subsec:holografischedisplay}
Eine der vielversprechendsten Forschungsrichtungen ist die Entwicklung holographischer \ac{NED}s auf Basis von Metaoberfläche (Metasurfaces) - planare Nanostrukturen, die Licht auf Subwellenlängen-Skala präzise manipulieren können. Bisherige \ac{NED}-Systeme erzeugen zweidimensionale Bilder auf einer festen optischen Tiefenebene; holografische Systeme hingegen rekonstruieren vollständige Wellenfronten und ermöglichen damit die Darstellung physikalisch korrekter dreidimensionaler Bilder mit natürlichen Tiefencues. Ein wegweisendes Ergebnis wurde 2024 in \textit{Nature} veröffentlicht: Forscher der Stanford University und der University of Washington demonstrierten ein vollfarb-3D-holografisches \ac{AR}-Display, das eine linsenlose holografische Lichtmaschine mit einer Metaoberflächen-Wellenleiteroptik kombiniert und dabei \ac{KI}-gestützte Algorithmen zur Optimierung der Wellenfrontberechnung einsetzt - mit dem Ziel, erstmals kompakte 3D-Holografie in einem brillenartigen Formfaktor zu realisieren. Ergänzend demonstrierten Forscher 2024 die Integration von Metaoberflächen in Kontaktlinsen, die holografische Inhalte direkt auf die Netzhaut projezieren - ein Ansatz, der den Formfaktor eines \ac{AR}-Systems auf ein Minimum reduziert.

\subsection{Micro-LED-Displays}
\label{subsec:MicroLed}
Micro-LEDs gelten als Schlüsseltechnologie der nächsten Displaygeneration für \ac{AR}-Anwendungen, da sie gegenüber OLEDs überlegene Luminanz, Lebensdauer und Energieeffizienz bieten. Das zentrale ungelöste Problem war bisher die Effizienzreduktion roter Micro-LEDs bei abnehmender Pixelgröße sowie die fehleranfällige Massentransferfertigung. Im Januar 2026 präsentierten Forscher des \ac{KAIST} einen Durchbruch: Durch eine AlInP/GaInP-Quantentopfstruktur und monollithische 3D-Integration konnte ein rotes Micro-LED-Display mit einer Pixeldichte von 1.700PPI realisiert werden - drei- bis viermal höher als aktuelle Smartphone-Displays. Parallel dazu entwickelten \ac{CEA}-Leti und das \ac{CRHEA} 2025 ein Verfahren zur Erzeugung von nativen RGB-Licht aus einem einzigen InGaN-Materialsystem auf Nanopyramiden-Basis, was die bisherige Notwendigkeit der Kombination verschiedener Materialsysteme für die drei Grundfarben in einem gemeinsamen Pixel beseitigt und einen Pfad zu vollintegrierten Vollfrag-Micro-LED-Mikrodisplays mit Pixelpitches unter 10 Mikrometern eröffnet.

\subsection{Eye-Tracking und Foveated Rendering}
\label{subsec:eyetrackundfoveated}
Eine weitere zentrale Forschungsrichtung verbindet die Hardware-Ebene der \ac{AR}-Brille mit der Rendering-Pipeline: Foveated Rendering nutzt Echtzeit-Eye-Tracking, um Bildinhalte nur im Bereich der Fovea - der zentralen, hochauflösenden Zone der Netzhaut - in voller Quelität zu rendern, während periphere Bereiche mit reduzierter Auflösung berechnet werden. Da die Fove nur ca. 2\textdegree des Sehfelds abdeckt, ermöglicht dieser Ansatz erhebliche Einsparungen  bei Rechenleistung und Energieverbrauch ohne wahrnehmbare Qualitätseinbuße. Aktuelle Forschungen adressieren dabei die Herausforderung der Prädiktionslatenz: Das System muss die Blickrichtung voraussagen, bevor das Auge tatsächlich eine Sakkade vollführt, da selbst kurze Verzögerungen zwischen Blickbewegung und Renderpunkt wahrnehmbar sind. Neuste laser-interferometrische Eye-Tracking-Sensoren erreichen dabei eine Genauigkeit von unter 1,62\textdegree bei einem Bruchteil des Energieverbrauchs konventioneller kamerabasierter Systeme, was sie für den mobilen Einsatz in \ac{AR}-Brillen besonders attraktiv macht.

\subsection{KI-gestützte Rendering- und Kontextverarbeitung}
\label{subsec:KIgestuetztRenderingundKontextverarb}
Die Integration von Methoden des maschinellen Lernens in die \ac{AR}-Pipeline gewinnt als Forschungsrichtung erheblich an Bedeutung. Auf der Rendering-Ebene ermöglichen neuronale Rendering-Verfahren die \ac{KI}-gestützte Approximation physikalischer Lichtausbreitung, was rechenintensive klassicshe Raytracing-Verfahren ersetzten oder ergänzen kann. Darüber hinaus werden \ac{KI}-Modelle eingesetzt, um die räumliche Umgebung in Echtzeit semantisch zu verstehen - Objekte werden nicht nur zu erkenn, sondern deren Kontext zu interpretieren - und damit kontextsensitive \ac{AR}-Überlagerungen zu erzeugen, die über statische Informationsanzeige hinausgehen. Die Konvergenz von Edge-Computing-Architekturen, Sensorfusion und großen vortrainierten Sprachmodellen ermöglicht nach aktuellen Forschungsprogrnosen \ac{AR}-Systeme, die als kontinuierlich mitdenkende Assistenzsysteme agieren, anstatt lediglich vordefinierte Inhalte darzustellen